% concatenated from paper.tex
\documentclass[11pt,a4paper]{article}

\usepackage{abstract}
\usepackage{amsmath}
\usepackage{amssymb}
\usepackage{amsthm}
\usepackage{hyperref}
\usepackage{booktabs} % For pretty tables
\usepackage{caption} % For caption spacing
%\captionsetup{aboveskip=1pt, belowskip=1pt} % Adjust as needed

\usepackage[a4paper,
  left=3.5cm,
  right=3.5cm,
  top=3cm,
  bottom=3cm
]{geometry}

\usepackage{subcaption} % For sub-figures
\usepackage{graphicx}
\usepackage{pgfplots}
\usepackage{adjustbox} % Add in the preamble
\usepackage[ruled,vlined]{algorithm2e}
%\usepackage{algpseudocode,algorithm,algorithmicx}

\usepackage[all]{nowidow}
\usepackage[utf8]{inputenc}
\usepackage{tikz}
\usetikzlibrary{er,positioning,bayesnet}
\usepackage{multicol}
\usepackage{minted}
\usepackage[inline]{enumitem} % Horizontal lists
% Used for displaying a sample figure. If possible, figure files should
% be included in EPS format.
%
% If you use the hyperref package, please uncomment the following line
% to display URLs in blue roman font according to Springer's eBook style:
% \renewcommand\UrlFont{\color{blue}\rmfamily}

\hyphenation{Bayes-ian}
\newcommand{\card}[1]{\left\vert{#1}\right\vert}
\newcommand*\Let[2]{\State #1 $\gets$ #2}
\definecolor{blue}{HTML}{1F77B4}
\definecolor{orange}{HTML}{FF7F0E}
\definecolor{green}{HTML}{2CA02C}

\pgfplotsset{compat=1.14}

\renewcommand{\topfraction}{0.85}
\renewcommand{\bottomfraction}{0.85}
\renewcommand{\textfraction}{0.15}
\renewcommand{\floatpagefraction}{0.8}
\renewcommand{\textfraction}{0.1}
\setlength{\floatsep}{3pt plus 1pt minus 1pt}
\setlength{\textfloatsep}{3pt plus 1pt minus 1pt}
\setlength{\intextsep}{3pt plus 1pt minus 1pt}
\setlength{\abovecaptionskip}{2pt plus 1pt minus 1pt}
%========================================

%%%%
\newcommand{\g}{{\, |\,}}
\newcommand{\Pp}{{\mathbb P}}
\newcommand{\wt}{{\widetilde t}}
\newcommand{\kauno}{{\mathcal L}}
\newcommand{\Nn}{{\mathcal N}}
\newcommand{\Tt}{{\Gamma}}
\newcommand{\oT}{\boldsymbol{T}}
\newcommand{\fa}{\mathbb{\forall}}
\newcommand{\upers}{\mathbb{\supset}}
\newcommand{\N}{\mathbb{N}}
\newcommand{\F}{\mathcal{F}}
\newcommand{\B}{\mathcal{B}}
\newcommand{\Rt}{\mathbb{R}^3}
\newcommand{\upos}{\subset}
\newcommand{\D}{\Delta}
\newcommand{\Ra}{\Rightarrow} % =>
\newcommand{\lra}{\longrightarrow} % ---> 
\newcommand{\subspace}{\hookrightarrow}
\newcommand{\R}{\mathbb{R}}
\newcommand{\Z}{\mathbb{Z}}
\newcommand{\U}{\mathcal{U}}
\newcommand{\C}{\mathbb{C}}
\newcommand{\tomh}{\cap}
\newcommand{\qq}{\qquad}
\newcommand{\q}{\quad}
\newcommand{\bd}{\partial}
\newcommand{\cl}{\overline}
\newcommand{\ektos}{\setminus}
\newcommand{\weakto}{\rightharpoonup}
\DeclareMathOperator*{\argmax}{argmax}
\DeclareMathOperator*{\argmin}{argmin}
%================================

%=================
\newtheorem{theorem}{Theorem}
\newtheorem{definition}{Definition}
\newtheorem{proposition}{Proposition}
\newtheorem{lemma}{Lemma}
\newtheorem{remark}{Remark}

%================================
\begin{document}
%
\title{Random tree Besov priors: Data-driven regularisation parameter selection}

%\author{ Hanne Kekkonen and Andreas Tataris}
% \author{%
%   Hanne Kekkonen\thanks{Email: \texttt{h.n.kekkonen@tudelft.nl}}\\
%   Delft University of Technology
%   \and
%   Andreas Tataris\\
%   Delft University of Technology
% }

\author{%
Hanne Kekkonen\thanks{Email: \texttt{h.n.kekkonen@tudelft.nl}}\hspace{1mm} and Andreas Tataris\\[0.6em]
\small{Delft University of Technology}\\
\small{Delft Institute of Applied Mathematics}\\
\small{Delft, The Netherlands}
}


%
\maketitle              % typeset the header of the contribution
%
%\begin{abstract}
%Notes on the EFWI and the Marchenko %integral equations.
%\keywords{}
%\end{abstract}
%
%% 
%$$
\begin{abstract}

We develop a data-driven algorithm for automatically selecting the regularisation parameter in Bayesian inversion under random tree Besov priors.
One of the key challenges in Bayesian inversion is the construction of priors that are both expressive and computationally feasible. Random tree Besov priors, introduced in \cite{Kekkonen2023}, provide a flexible framework for capturing local regularity properties and sparsity patterns in a wavelet basis. In this paper, we extend this approach by introducing a hierarchical model that enables data-driven selection of the wavelet density parameter, allowing the regularisation strength to adapt across scales while retaining computational efficiency. We focus on nonparametric regression and also present preliminary plug-and-play results for a deconvolution problem.\\

\noindent \textit{Keywords:} Bayesian inverse problems, statistical inversion, hierarchical Bayes, Besov priors, wavelets, multiscale sparsity.

\end{abstract}
\section{Introduction}

Inverse problems arise from the need to estimate unknown quantities from indirect, often incomplete, noisy measurements. They are encountered in various applications, ranging from medical imaging and remote sensing to geophysical exploration and financial modelling. Inverse problems are typically ill-posed, meaning that small perturbations in data or inaccuracies in the model can result in large errors or instabilities in the reconstructions. To obtain stable reconstructions, we must carefully combine prior knowledge about the unknown with the measurements.

In practice, measurements are always finite-dimensional and contaminated by noise. To avoid undesirable artefacts and potential changes in the qualitative behaviour of reconstructions due to discretisation, it is advisable to construct a continuous model that is discretised as late as possible in the analysis and solution process \cite{Haario2004, Lassas2004, Stuart2010}. This leads to a continuous model of the form
\begin{equation*}\label{eq:IP_continuous}
M = Af + \varepsilon \mathbb{W},
\end{equation*}
where $A:X\to Y$ is a linear operator between Banach spaces $X$ and $Y$, $\varepsilon>0$ describes the noise amplitude, and $\mathbb{W}$ is a Gaussian white noise process indexed by $Y$ \cite{Reiss2008}. 

Bayesian inversion provides a powerful framework for addressing inverse problems by combining the indirect, noisy data with prior information in a probabilistic manner. This offers an attractive alternative to regularisation by also delivering information on how uncertainties in the data or forward model affect the solution.
The choice of the prior distribution, which encodes assumptions about the unknown quantity, is central to the Bayesian approach. While Gaussian priors offer computational efficiency and theoretical simplicity, they impose global smoothness constraints, making them less effective for problems such as image and signal reconstruction, where capturing sharp edges and interfaces is essential.

Recently, Besov priors have emerged as a compelling alternative, bridging the gap between rigorous mathematical formulation and practical performance \cite{Lassas2009, Dashti2012}. Besov priors naturally capture local smoothness properties while allowing discontinuities, making them well-suited for problems involving spatially varying regularity \cite{Bui2015, Cui2021, Hamalainen2013, Horst2025, Kolehmainen2012, Rantala2006, Suuronen2020, Wang2017}. 

However, a limitation of traditional Besov priors constructed via the Karhunen-- Loève expansion is their inability to explicitly encode complex covariance structures. To overcome this limitation, random tree Besov priors were introduced in \cite{Kekkonen2023}. Incorporating a hierarchical tree structure into the wavelet representation allows for more flexible modelling of local variations and rapid spatial changes. Although this approach does not directly encode a covariance structure comparable to Gaussian processes, it significantly enhances flexibility compared to classical Besov priors. The random tree Besov priors have been applied, for example, to elliptic PDEs \cite{Schwab2024, Stein2023}. 

Random tree Besov priors are conceptually related to other tree-structured Bayesian models. In finite-dimensional signal and image processing, wavelet-domain hidden Markov tree models \cite{Crouse2002,Romberg2001} represent wavelet coefficients via Gaussian mixtures with hidden states forming a Markov tree, capturing clustering of large and small coefficients and cross-scale persistence. In nonparametric density estimation tree-based priors such as Pólya trees\cite{Castillo2017, Castillo2021, Castillo2022}, as well as Bayesian CART-type priors with multiscale shrinkage \cite{Castillo2019}, also exploit hierarchical tree structures. In contrast, random tree Besov priors are constructed directly as Besov function-space priors in a wavelet basis and are tailored for Bayesian inverse problems.

The construction proposed in \cite{Kekkonen2023} leveraged wavelet decompositions, modelling $f$ as a random function whose wavelet coefficients are supported on structured trees. This tree-pruning approach naturally promotes sparsity and local adaptivity by enforcing hierarchical constraints for the wavelet coefficients. The hierarchical tree structure ensures that a wavelet coefficient is selected only if all of its ancestors in the decomposition are also selected. This condition is motivated by the observation that true discontinuities in a signal or image produce large wavelet coefficients consistently across multiple scales, from coarser to finer resolutions. In contrast, noise-induced coefficients typically appear large only at the finest scales and rapidly diminish at coarser levels due to the inherent averaging effect of wavelets.

In \cite{Kekkonen2023} the wavelet density index $\beta$ that controls the sparsity of the wavelet coefficients was selected manually. Although this method yielded promising results, the optimal choice of $\beta$ remained challenging and data-dependent. Furthermore, employing a single fixed regularisation parameter across all wavelet decomposition levels is overly restrictive, since wavelet coefficients typically exhibit varying degrees of sparsity at different resolutions: fine scales often have sparser representations than coarser scales \cite{Adcock2017}.

In this paper, we address these limitations by presenting an automatic data-driven method for selecting the wavelet density index $\beta$. We introduce a hierarchical Bayesian model in which the wavelet density parameter is allowed to vary across decomposition levels, resulting in a vector-valued parameter $\boldsymbol{\beta}= (\beta_j)_j$. We place a suitable hyperprior on $\boldsymbol{\beta}$ and propose a computationally efficient method for estimating the optimal $\boldsymbol{\beta}$ directly from data by considering the levelwise maximum a posteriori (MAP) estimates. This strategy allows the regularisation strength to adapt flexibly to the scale-specific sparsity inherent in wavelet decompositions and removes manual tuning of $\beta$. We introduce an efficient recursive algorithm for computing the MAP estimates of the wavelet coefficients, tree structure, and regularisation parameters simultaneously for both Gaussian and Laplace base prior distributions on the wavelet coefficients. We provide numerical experiments demonstrating that our automatically chosen $\boldsymbol{\beta}$ performs at least as well as the optimally chosen fixed $\beta$. 

The remainder of this paper is organised as follows. In Section \ref{Sec:Preliminaries}, we review relevant concepts for Besov spaces and wavelet-based random variables, and define the random tree Besov priors in detail. Section \ref{Sec:BetaSelection} describes our proposed Bayesian inference method for automatically selecting level-dependent regularisation parameters $\boldsymbol{\beta}$. Numerical examples illustrating the performance and practical advantages of our approach are provided in Section \ref{Sec:Numerical}. 



% contributions
\section{Random tree Besov priors\label{Sec:Preliminaries}}
%\subsection{Bayesian Inversion}

We briefly recall the essential theory of Besov spaces and the construction of Besov measures in the next subsection, before moving to the hierarchical random tree structure.



\subsection{Besov random variables}\label{Sec:BesovMeasures}
In this paper we use the wavelet representation of Besov norms. 
Let $\phi$ and $\Psi^\ell$, $\ell=1,\ldots ,2^d-1$, be compactly supported scaling and wavelet functions suitable for multi-resolution analysis of smoothness $C^r$ with some $r\in\N_+$. We denote
\begin{align*}
\psi^\ell_{jk}(x)&=2^{jd/2}\Psi^\ell(2^jx-k),\qquad k\in\Z^d,\;  
j\geq 0,  \; \ell\in\kauno:=\{1,\ldots, 2^d-1\},\\
\phi_{k}(x)&= \phi(x-k), \qquad k\in \Z^d.
\end{align*}
Any function $f\in L^2(\R^d)$ can then be decomposed in terms of the wavelet basis
\begin{align*}
f(x)&=\sum_{k\in\Z^d}\langle f, \phi_k\rangle\phi_k(x)
+\sum_{\substack{j\geq 0,\ k\in \Z^d\\ \ell\in\kauno}}\langle f, \psi^\ell_{ jk}\rangle\psi^\ell_{ jk}(x)\\
&=\sum_{k\in\Z^d}f_{-1k}\phi_k(x)+\sum_{\substack{j\geq 0,\ k\in \Z^d
\\ \ell\in\kauno}}f^\ell_{ jk}\psi^\ell_{ jk}(x).
\end{align*}

The Besov space $B_{pp}^s(\R^d)$ is defined as
\begin{align*}
    B_{pp}^s(\R^d):= \{f\in L^2(\R^d): \|f\|_{B^s_{pp}(\R^d)}<\infty \},
\end{align*}
with the norm
\begin{align*}
\|f\|_{B^s_{pp}(\R^d)}
= \left(\sum_{{j=-1}}^\infty 2^{jp\big(s+\frac d2-\frac dp\big)}
\|\textbf{f}_j\|_{\ell^p}^p\right)^{1/p},
\end{align*}
where
\begin{align*}
\textbf{f}_j&=(f^\ell_{ jk})_{k\in \Z^d, \; \ell\in\kauno}\quad \text{with}\quad
\|\textbf{f}_j\|_{\ell^p}=\left(\sum_{k\in \Z^d, \; \ell\in\kauno}
|f^\ell_{ jk}|^p\right)^{1/p}  \qquad \textrm {if}\;\, j\geq 0,\quad \textrm{and}\\
\textbf{f}_{-1}&=(f_{-1k})_{k\in \Z^d} \quad \text{with}\quad
\|\textbf{f}_{-1}\|_{\ell^p}=\left(\sum_{k\in \Z^d}|f_{-1k}|^p\right)^{1/p}.
\end{align*}
For more details on Besov spaces we refer to 
\cite{Adams1975,Gine2015,Tartar2007}.

We note that Besov spaces are suitable for representing functions with local smoothness variations, including sharp edges. For example, the Besov space $B^1_{11}$ is close to the space of total variations in the sense that 
\begin{align*}
    B^{1}_{11}(\R^d)\subset W^{1,1}_{loc}(\R^d)\subset B^{1-\varepsilon}_{11}(\R^d), \qquad \textrm{for all}\quad \varepsilon >0.
\end{align*}

For simplicity, we ignore boundary effects below by considering functions with support in the interior of the unit cube $D=[0,1]^d$. 
We could also adjust the definition of the Besov priors for different boundary conditions, and study the behaviour  near the boundary, but since this is application-specific we will not consider it in this paper.

Let $(X^\ell_{jk})_{\ell\in\kauno,\, j\in\N,\,k\in\Z^d}$ be an i.i.d. sequence of real random variables with probability density function $\pi(x)\propto\exp (-\frac{1}{2\kappa^p}|x|^{p})$, $1\leq p<\infty$, denoted by $X^\ell_{jk}\sim \Nn_p(0,\kappa^p)$.
Define the random function
\begin{align}\label{eq:BesovPrior}
f(x)=\sum_{\substack{\ell\in\kauno,\, j\in\N\\ k\in\Z^d}}
2^{-j\big(s+\frac d2-\frac dp\big)} X^\ell_{ j,k}\psi^\ell_{j,k}(x),\quad x\in D,
\end{align}
with $s<r$.  We call $f$ a $B^{s}_{pp}$-random variable.

These random variables take values almost surely in Besov spaces $B^{t}_{pp}$ with $t<s-\frac{d}{p}$, while a realisation $f$ belongs to $B^{s-d/p}_{pp}$ with probability zero \cite{Dashti2012}. Informally, $f$ has density proportional to $\exp(-\frac{1}{2\kappa^p}\|f\|_{B^s_{pp}}^p)$, see \cite{Dashti2012,Lassas2009}. 

If $p=2$, \eqref{eq:BesovPrior} defines a Gaussian measure with Cameron-Martin space $B^s_{22}=H^s$. For Gaussian measures, the Cameron–Martin space characterises admissible translations under which the measure remains equivalent, and governs small-ball probabilities, concentration, and posterior contraction. In contrast, Besov priors require a more delicate analysis. For $p\in[1,2]$, the Hilbert space $B^s_{22}$ describes admissible translations, while the Banach space $B^s_{pp}$ determines concentration and contraction rates \cite{Agapiou2021}. 

Since the Besov priors are constructed using a Karhunen--Lo\`eve expansion they  lack a natural covariance structure, limiting control over local correlation properties. To overcome this shortcoming random tree Besov priors introduce hierarchical wavelet coefficient selection to enhance flexibility in representing rapid local variations, mimicking the capabilities typically associated with covariance structures in Gaussian processes.\\


\subsection{Random tree Besov variables}\label{Sec:RandomTreeBesovVariables}


We will now introduce the random tree Besov priors.
Since we consider functions in the interior of $D=[0,1]^d$ we see that
$f^\ell_{ jk}=0$ if $k_l<0$ or $k_l>2^j-1$ for some $l=1,\dots,d$. 
We call the remaining set of indices of the wavelet coefficients the {\it entire tree} and denote this by
\begin{align*}
\oT=\{(j,k)\in \N\times \N^d\,,\, j\geq 0,\ k=(k_1,\dots,k_d),\
0\leq k_l\leq 2^j-1\}.
\end{align*}

Note that the index $\ell$ in Section \ref{Sec:BesovMeasures} denotes the number of parallel trees. For $1$D signals $\ell=1$ and we can organise the wavelet coefficients into a binary tree. For $2$D images $\ell\in{1,2,3}$ making the structure more complex since we need to consider three parallel quadtrees: one for vertical, horizontal, and diagonal subbands. 

The idea of the random tree Besov priors presented in \cite{Kekkonen2023} was to use the Karhunen-Lo\`eve expansion with a wavelet basis to create priors similar to \eqref{eq:BesovPrior} but to choose non-zero wavelets in the sum in a systematic way. To achieve this we will introduce the set of {\it proper subtrees}:
\begin{align*}
\Gamma=\{T\subset\oT\,|\,
\hbox{if $(j,k)\in T$ then $(j-1,[k/2])\in T$}\},
\end{align*}
where $[k/2]=([k_1/2],\dots,[k_d/2])$ is the vector whose
elements are the integer parts of $k_l/2$. This means that for any $T\in\Gamma$ and any node $(j,k)\in T$, all of its ancestors also belong to the tree $T$, and therefore every node in a proper subtree is connected to the root node. 
The random subtree can be generated by a Galton-Watson-type branching mechanism on the dyadic tree, where each potential child node is included independently with probability $\beta$ conditional on its parent being included.

For simplicity,  we assume below that all the bands $\ell$ have the same subtree structure. This assumption can easily be relaxed by considering separate subtrees $T^\ell$ instead.  
We can now introduce the random tree Besov priors.

\begin{definition}\label{Def:TreeRandomVariable}
Let $\{\psi^\ell_{jk}\}_{\ell\in\kauno, j\in\N,\,k\in\Z^d}$ be an $r$-regular wavelet basis for $L^2(\R^d)$, and $\beta=2^{\gamma-d}$ with some $\gamma\in(-\infty,d\,]$. Consider pairs $(X,T)$ where $X$ is an $\R^{\oT^{2^d-1}}$-valued random variable, and $T\in {\Tt}$ is a random tree. We assume that $X$ and $T$ are independent random variables, having the following distributions
\begin{itemize}
\item The sequence $X$ consists of i.i.d random variables  with probability density proportional to $\exp (-\frac{1}{2\kappa^p}|x|^{p})$, $\kappa>0$ and $1\leq p <\infty$, that is, $X_{jk}^\ell\sim\Nn_p(0,\kappa^p)$.
\item The proper subtree $T$ is built recursively. The root node is always included. For each other node, conditional on its parent being in $T$, the node is included independently with probability $\beta$; if its parent is not in $T$, the node is excluded.
\end{itemize}
Let $f$ be the random function
\begin{align*}
f(x)=\sum_{\substack{(j,k)\in T\\ \ell\in\kauno}}
2^{-j\big(s+\frac d2-\frac dp\big)}X^\ell_{j,k}\psi^\ell_{j,k}(x),\quad x\in D,
\end{align*}
where $s<r$. We say that $f$ is a $B^{s}_{pp}$-random variable with wavelet density index $\beta$.
\end{definition}

In \cite{Kekkonen2023} the wavelet density index was chosen manually but we would like to choose $\beta$ automatically from data.
Below we will consider level dependent $\boldsymbol{\beta}=(\beta_j)_j$ to allow varying sparsity levels since it is common that the wavelet coefficients of a signal or image are sparser at the finer resolution scales than at the coarser scales. 
We will assign a hyperprior for $\boldsymbol{\beta}$ and consider the levelwise maximum a posteriori (MAP) estimates for the wavelet density index. 


\section{Tree pruning with automatic $\boldsymbol{\beta}$ selection \label{Sec:BetaSelection}}

In this section we develop our tree–pruning scheme with automatic selection of the level–dependent wavelet density indices $\boldsymbol{\beta}=(\beta_j)_{j=1}^J$. We begin by introducing the problem and recall the tree–pruning algorithm for a fixed, known $\beta$ in Section \ref{Sec:PruningFixedBeta} in order to set notation and highlight the role of $\beta$ as a regularisation parameter. Sections \ref{Sec:PruningGaussian} and \ref{Sec:PruningLaplace} then describe our main contribution: a hierarchical Bayesian model in which each $\beta_j$ is equipped with a hyperprior and is estimated from the data via levelwise MAP optimisation, for the Gaussian and Laplace base priors respectively. 

We consider the nonparametric regression problem where
we have observations at $2^{d(J+1)}$, $J\in\N$, regularly spaced points $x_i$ on $D=[0,1]^d$, $d=1,2$, of some unknown function $f$ subject to noise
\begin{align}\label{eq:DenoisingMain}
    M_i=f(x_i)+W_i,
\end{align}
where $W_i\sim \Nn(0,\sigma^2)$ is white noise independent of $f$. We address the cases $p=1$ and $p=2$ which correspond to the non-zero wavelet coefficients of $f$ being drawn from the Laplace and Gaussian distributions respectively. 

The wavelet coefficients can be arranged in a finite tree of the form
\begin{equation*}
\oT=\{(j,k)\in \N\times \N^d\ |\ 0\le j\le J,\ k=(k_1,\dots,k_d),\ 0\leq k_l\leq 2^j-1\}
\end{equation*}
and the non-zero wavelet coefficients will be in a finite proper subtree
\begin{align*}
\Gamma=\{T\subset\oT\,|\,
\hbox{if $(j,k)\in T$ then $(j-1,[k/2])\in T$}\},
\end{align*}
where $[k/2]$ is the integer part of $k/2$. 
We denote $(j',k')\prec(j,k)$ if $(j,k)$ is an ancestor of $(j',k')$.
The complete subtree with a root node $(j,k)$ can then be written as
\begin{align*}
\oT_{(j,k)} =\{(j',k')\in \oT\ |\ (j',k')\preceq (j,k)\}.
\end{align*}

We model the unknown wavelet density indices $\beta_j$, $j=1,\dots,J$, with a hyperprior
\begin{align*}
\pi(\beta_j)=2^{(1+a)}(1+a)(0.5-\beta_j)^a, 
%\pi(\beta_j)\propto(0.5-\beta_j)^a 
\quad \beta_j\in[0,0.5],\ a\geq0.  
\end{align*}
We restrict $\beta_j\in [0,0.5]$ so that the prior acts as a sparsity-promoting regulariser. For $\beta_j<0.5$ we have $-\log(\beta_j)>-\log(1-\beta_j)$, so adding nodes is discouraged \textit{a priori}. The strength of the hyperprior is controlled by the parameter $a$: as $a$ increases, the prior penalises large $\beta_j$ more heavily, leading to sparser trees and consequently stronger regularisation. 


To generate a random subtree with wavelet density $\beta=(\beta_j)_{j=1}^J$, we first consider the nodes independently and set the probability for a node to be included as $\Pp(t_{jk}=1 \g \beta_j)=\beta_j$, and correspondingly the probability for exclusion as $\Pp(t_{jk}=0 \g \beta_j)=1-\beta_j$. 
The resulting proper subtree $T$ then consists of all nodes that connect to the root node, that is, the nodes $(j,k)$ for which
\begin{align*}
\wt_{jk}=\prod_{(j',k')\succeq(j,k)}t_{j'k'} =1.
\end{align*}

We model the unknown function $f$ using a truncated random tree Besov prior and write 
\begin{align}\label{eq:DefPrior}
\begin{split}
f(x) & =\sum_{(j,k)\in T} \langle f,\psi_{jk}\rangle\psi_{jk}(x)\\
& =\sum_{(j,k)\in \oT} \wt_{jk} g_{jk}\psi_{jk}(x),
\end{split}
\end{align}
where $g_{jk} \sim \Nn(0,\kappa^2)$ or $g_{jk} \sim \text{Laplace}(0,\kappa)$, and $\wt_{jk}\in\{0,1\}$ defines if a node $(j,k)\in \oT$ is included as described above. The variables $g_{jk}, t_{jk}$ are assumed to be mutually independent.

We consider noise of the form $W=\sum w_{jk}\psi_{jk}$ where $w_{jk}\sim \Nn(0,\sigma^2)$. 
The data can then be written as $M=\sum m_{jk}\psi_{jk}(x)$, where
\begin{align*}
m_{jk}=\left\{\begin{array}{ll}g_{jk}+w_{jk}, & \textrm{when}\ \wt_{jk}=1\\
w_{jk}, & \textrm{when}\ \wt_{jk}=0. \end{array} \right.
\end{align*}

Since $g_{jk}, t_{jk}$ are assumed to be mutually independent, the posterior distribution of $\boldsymbol{g}=(g_{jk})_{(j,k)\in\oT}$, $\boldsymbol{t}=(t_{jk})_{(j,k)\in\oT}$, and $\boldsymbol{\beta}=(\beta_j)_{j=1}^J$, given the data $\boldsymbol{m}=(m_{jk})_{(j,k)\in\oT}$, can be written as
\begin{align*}
\pi(\boldsymbol{g},\boldsymbol{t},\boldsymbol{\beta}\ |\ \boldsymbol{m}) & \propto\pi(\boldsymbol{m}\ |\ \boldsymbol{g},\boldsymbol{t})\pi(\boldsymbol{g})\pi(\boldsymbol{t}\g \boldsymbol{\beta})\pi(\boldsymbol{\beta})\\
& =\prod_{(j,k)\in \oT} \pi(m_{jk}\ |\ g_{jk},\wt_{jk})\pi(g_{jk})\pi(t_{jk}\g\beta_j)\pi(\beta_j)\\
& = \prod_{(j,k)\in \oT} z_{jk},
\end{align*}
where we have denoted $z_{jk}=\pi(g_{jk},\widetilde{t}_{jk},\beta_j \g m_{jk})$.

The posterior for $g_{jk}$ and $\beta_j$ when the node $(j,k)$ is included in the tree is given by
\begin{align*}
z_{jk}^1 & =\pi(g_{jk},\beta_j \g m_{jk},\widetilde{t}_{jk}=1)\\
%=\pi(m_{jk}\ |\ g_{jk},1)\pi(g_{jk})\pi(1 \g \beta_j)\pi(\beta)\\
& \propto \exp\Big(-\frac{1}{2\sigma^2}(m_{jk}-g_{jk})^2-R(g_{jk})\Big)\beta_j\Big(\frac{1}{2}-\beta_j\Big)^a
\end{align*}
and the posterior when the node is excluded is 
\begin{align*}
z_{jk}^0 & =\pi(g_{jk},\beta_j \g m_{jk},\widetilde{t}_{jk}=0)\\
%=\pi(m_{jk}\|\ g_{jk},0)\pi(g_{jk})\pi(0\g \beta_j)\pi(\beta)\\
& \propto \exp\Big(-\frac{1}{2\sigma^2}m_{jk}^2-R(g_{jk})\Big)(1-\beta_j)\Big(\frac{1}{2}-\beta_j\Big)^a, 
\end{align*}
where $\pi(g_{jk})\propto\exp(-R(g_{jk})),$ with $R(g_{jk})={g_{jk}^2}/(2\kappa^2)$ when $g_{jk}\sim\Nn(0,\kappa^2)$ and $R(g_{jk})={|g_{jk}|}/\kappa$ when $g_{jk}\sim\text{Laplace}(0,\kappa)$. 


\subsection{Tree pruning with a fixed $\beta$}\label{Sec:PruningFixedBeta}

The case in which the wavelet density index $\beta$ is fixed and given was considered in \cite{Kekkonen2023}. We will discuss this algorithm shortly for completeness. For simplicity, we will describe the case $d=1$ below. 

If $\beta$ is fixed, we can write the optimisation problem in the form
\begin{align}\label{eq:Optimisegt}
\begin{split}
(\hat{\boldsymbol{g}},\hat{\boldsymbol{t}}) & = \argmax_{\boldsymbol{g},\boldsymbol{t}} \pi(\boldsymbol{g},\boldsymbol{t}\g \boldsymbol{m})\\
& = \argmax_{\boldsymbol{g},\boldsymbol{t}} \Bigg[ z^1_{00}
 +\sum_{k=0}^1 \bigg( t_{1k}\Big(z^1_{1k}
 + \!\!\! \sum_{(j',k')\prec(1,k)} \!\!\! z_{j'k'}\Big)
 + (1-t_{1k}) \!\!\! \sum_{(j',k')\preceq(1,k)} \!\!\! z^0_{j'k'}\bigg)\Bigg]
\end{split}
\end{align}  
where we have assumed that the root node $(0,0)$ is always chosen. Above the optimisation problem is split into two subproblems concerning optimisation of the left and right subtrees with root nodes $(1,0)$ and $(1,1)$. A new root node $(1,k)$, $k=0,1$, is either included in the tree or excluded. Including the root node leads to two new optimisation problems, while excluding the node prunes the whole branch. Since we do not know the weight of the optimised subtree we cannot yet compare whether including the node or pruning the branch is an optimal choice. To solve the optimisation problem we continue the recursive splitting until the bottom level $j=J$, where $\wt_{Jk}=t_{Jk}$, is reached. At the bottom level we simply compare $z^1_{Jk}$ and $z^0_{Jk}$ componentwise. The optimal tree can then be built from bottom up by comparing the weight of an optimised subtree to the weight of pruning the branch. 

In \cite{Kekkonen2023} solving the MAP estimate for $\boldsymbol{g}$ and $\boldsymbol{t}$ with a fixed $\beta$ was shown to lead to a simple tree pruning algorithm when $g_{jk}\sim\Nn(0,1)$ and to tree pruning with additional soft thresholding when $g_{jk}\sim \text{Laplace}(0,\kappa)$. In tree pruning, all branches that do not carry enough information are pruned, resulting in zero wavelet coefficients in the sum \eqref{eq:DefPrior}. How much information is considered enough is determined by the wavelet density index $\beta$ which acts as a regularisation parameter. 

In the following we choose $\hat{\boldsymbol{\beta}}=(\hat{\beta})_{j=1}^J$ to be the levelwise MAP estimates and then use the selected $\hat{\beta}_j$ in the optimisation problem \eqref{eq:Optimisegt}. Similarly to $\hat{\boldsymbol{g}}$ and $\hat{\boldsymbol{t}}$ the wavelet density index $\hat{\boldsymbol{\beta}}$ is selected recursively starting from the  bottom level $j=J$, fixing the selection, and moving up a level.

Below we first consider the case $w_{jk}\sim \Nn(0,1)$ and $g_{jk}\sim\Nn(0,1)$. Proposition \ref{Prop:scaling} shows that in practice, under Gaussian base prior, the noise and prior variances $\sigma^2$ and $\kappa^2$ influence only the optimal choice of the wavelet density index $\beta$.

\subsection{Automatic $\beta$ selection with Gaussian draws}\label{Sec:PruningGaussian}

The problem of maximising the posterior $\pi(\boldsymbol{g},\boldsymbol{t},\boldsymbol{\beta} |\ \boldsymbol{m})$ is equivalent to minimising $-\log\big(\prod_{(j,k)\in \oT} z_{jk}\big)$.
On the bottom level $j=J$ a node is included in the tree if 
\begin{align*}
-\log(z^1_{Jk})\leq -\log(z^0_{Jk}),        
\end{align*}
that is, if the weight of including the node is smaller than the weight of excluding it. Notice that the penalty for including a node $-\log(\beta)$, compared to the penalty for excluding a node $-\log(1-\beta)$, is regularising only when $\beta<0.5$.

We assume Gaussian base prior $g_{jk}\sim\Nn(0,1)$ and noise $w_{jk}\sim \Nn(0,1)$. On the bottom level the weight resulting from including a node is
\begin{align*}
Y^1_{Jk}
= \min_{g_{Jk}}\big(-\log(z^1_{Jk})\big)
= \frac{1}{4}m^2_{Jk}-\log\bigg(\beta_J\Big(\frac{1}{2}-\beta_J\Big)^a\bigg)
\end{align*}
which is attained by $\hat{g}^1_{Jk}=m_{Jk}/2$. If a node is not included we set $\hat{g}^0_{Jk}=0$ leading to weight
\begin{align*}
Y^0_{Jk}
= \min_{g_{Jk}}\big(-\log(z^0_{Jk})\big)
= \frac{1}{2}m^2_{Jk}-\log\bigg((1-\beta_J)\Big(\frac{1}{2}-\beta_J\Big)^a\bigg).
\end{align*}

We include a node in the tree if selecting it leads to a lower cost than pruning it, that is, $Y^1_{Jk}\leq Y^0_{Jk}$ or equivalently
\begin{align*}
\frac{1}{4}m^2_{Jk}-\log(\beta_J)
& \leq \frac{1}{2}m^2_{Jk}-\log(1-\beta_J). 
\end{align*}
Notice that the prior for $\beta$ does not play a role in whether a node is included or not. 
We see that, for a fixed value of $\beta_J$, the node $(J,k)$ should be included if 
\begin{align}\label{eq:GaussianChosen}
m_{Jk}^2\geq 4\log\bigg(\frac{1-\beta_J}{\beta_J}\bigg).
\end{align}

Denote by $A_\beta$ the set of included nodes, that is, the values of $k$ for which \eqref{eq:GaussianChosen} holds for a fixed $\beta$. We want to find the wavelet density index $\hat{\beta}_J$ that minimises the weight of the bottom row 
\begin{align}
\begin{split}
B_{J,A_{\beta_J}}
& = \sum_{k\in A_{\beta_J}} Y^1_{Jk}
+\sum_{k\notin A_{\beta_J}} Y^0_{Jk}\\
& = \sum_{k\in A_{\beta_J}}\frac{1}{4}m^2_{Jk} 
+ \sum_{k\notin A_{\beta_J}}\frac{1}{2}m^2_{Jk}\\
& \quad\quad +|A_{\beta_J}|\log\Big(\frac{1-\beta_J}{\beta_J}\Big)
-n_J\log\bigg((1-\beta_J)\Big(\frac{1}{2}-\beta_J\Big)^a\bigg)
\end{split}
\end{align}
where $|A_{\beta_J}|$ is the number of included nodes and $n_J=2^J$ is the number of nodes on the bottom row. We denote the data fit term by
$$L(m_J,\beta_J)=\sum_{k\in A_{\beta_J}}\frac{m^2_{Jk}}{4}+\sum_{k\notin A_{\beta_J}}\frac{m^2_{Jk}}{2}$$
and the regularisation penalty by
$$R(\beta_J)=|A_{\beta_J}|\log\Big(\frac{1-\beta_J}{\beta_J}\Big)
-n_J\log\bigg((1-\beta_J)\Big(\frac{1}{2}-\beta_J\Big)^a\bigg).$$
and write 
\begin{align}\label{eq:WeightBottomGaussian}
B_{J,A_{\beta_J}} = L(m_J,\beta_J)+R(\beta_J).
\end{align}

As $\beta_J$ increases, more nodes are included in the tree. However, the set of included nodes changes only at the points 
\begin{align*}
\beta_{Jk} = \frac{1}{1+\exp\big(\frac{1}{4}m^2_{Jk}\big)}.
\end{align*}
Therefore, to minimise $B_{J,A_{\beta_J}}$, it is sufficient to consider the $2^J$ grid points $\beta_{Jk}$. 
The penalty term $R(\beta_J)$ is a convex function, while the data term $L(m_J,\beta_J)$ is a decreasing function of $\beta_J$ on the interval $[0,0.5]$. Therefore, there exists a unique minimiser $\hat{\beta}_J=\argmin_{\beta_{Jk}} B_{J,A_{\beta_{Jk}}}$. 
Note that the $\beta$ grid is data-determined, the hyperprior only affects the argmin over the grid.

In practice, we reorder the measurements $m_{Jk}$ on the bottom level in descending order and denote this permutation by $m_{J\sigma(k)}=m_{J(k)}$, so that
\begin{align*}
m_{J(1)}\geq \dots \geq m_{J(n_J)}. 
\end{align*}
This reordering corresponds to arranging the wavelet density index grid points from smallest to largest. For $\beta_{J(1)}$, only the node with the largest measurement $m_{J(1)}$ is chosen, while for $\beta_J=\beta_{J(n_J)}$ all of the nodes on the bottom row are selected. The levelwise optimal $\hat{\beta}_J$ is found as the $\beta_{J(k)}$ that results in the first negative difference $B_{J,A_{\beta_{J(k)}}}-B_{J,A_{\beta_{J(k+1)}}}$. 

Once the optimal $\hat{\beta}_J$ has been solved we can use it to calculate the optimal $g_{Jk}$ and $t_{Jk}$ as in \cite{Kekkonen2023}. After this we fix the values on the bottom row and move to the next level $j=J-1$. The process for the levels $j<J$ is similar to the bottom level but this time we compare the weight of an optimised subtree with root node $(j,k)$ to the weight of pruning the whole branch. 

Denote the weight of a previously optimised subtree with root node $(j,k)$ by $F_{jk}$. The weight of including the optimised subtree with a root node $(j,k)$ while paying the inclusion penalty is
\begin{align*}
Y^1_{jk}=F_{jk}-\log\bigg(\beta_j\Big(\frac{1}{2}-\beta_j\Big)^a\bigg),    
\end{align*}
while the weight of pruning the branch is 
\begin{align*}
Y^0_{jk} 
& = \frac{1}{2}\|m{|_{T_{jk}}}\|^2
-\log\bigg((1-\beta_j)\Big(\frac{1}{2}-\beta_j\Big)^a\bigg)
-\sum_{l=j+1}^J n_l \log\bigg((1-\hat{\beta}_l)\Big(\frac{1}{2}-\hat{\beta}_l\Big)^a\bigg)\\
& = \frac{1}{2}\|m{|_{T_{jk}}}\|^2
-\log\bigg((1-\beta_j)\Big(\frac{1}{2}-\beta_j\Big)^a\bigg)
-C_j,
\end{align*}
where $n_l=2^{l-j}$ is the number of the nodes in the subtree on the level $l$ and $\hat{\beta}_l$, $l=j+1,\dots,J$, are the previously optimised wavelet density indices. 

Notice that $C_j=\sum_{l=j+1}^J n_l \log((1-\hat{\beta}_l)(0.5-\hat{\beta}_l)^a)$ collects the penalties for not selecting a node on levels $l>j$ in the subtree $T_{jk}$. Since the corresponding wavelet density indices $\hat{\beta}_l$, $l>j$, have already been fixed in the bottom–up recursion, $C_j$ does not depend on $\beta_j$ and is a constant in the optimisation with respect to $\beta_j$.


As before we include the optimised subtree if it is cheaper than pruning the branch, that is, if
\begin{align*}
Y^1_{jk}\leq Y^0_{jk}.
\end{align*}
Solving the above we see that, given a measurement $m_{jk}$, a node $(j,k)$ is chosen if 
\begin{align*}
\beta_{jk}\geq \frac{1}{1+\exp(\frac{1}{2}\|m|_{T_{jk}}\|^2-F_{jk}-C_j)},
\end{align*}
where $F_{jk}$ is the weight of a previously optimised subtree with root node $(j,k)$. 
As on the bottom row the prior for $\beta$ does not play a role in the selection of the grid points $\beta_{jk}$. They depend only on the data. 

To find a levelwise optimal $\beta_j$ we want to minimise the weight the row adds to the tree. Namely, we want to find $\hat{\beta}_j$ that minimises
\begin{align*}
B_{j,A_{\beta_j}}
& = \sum_{k\in A_{\beta_j}} Y^1_{jk}
+\sum_{k\notin A_{\beta_j}} Y^0_{jk}\\
& = \sum_{k\in A_{\beta_j}}F_{jk} 
+ \sum_{k\notin A_{\beta_j}}\Big(\frac{1}{2}\|m|_{T_{jk}}\|^2-C_j\Big)\\
& \quad\quad +|A_{\beta_j}|\log\Big(\frac{1-\beta_j}{\beta_j}\Big)
-n_j\log\bigg((1-\beta_j)\Big(\frac{1}{2}-\beta_j\Big)^a\bigg)\\
& = L(m_j,\beta_j)+R(\beta_j)
\end{align*}
where the data term is now 
$$L(m_j,\beta_j)= \sum_{k\in A_{\beta_j}}F_{jk} 
+ \sum_{k\notin A_{\beta_j}}\Big(\frac{1}{2}\|m|_{T_{jk}}\|^2-C_j\Big).$$
We absorb the penalty term $C_j$ into the data term so that the cost of pruning a branch can be compared directly with the optimised subtree weight $F_{jk}$, which accounts for the penalties associated with optimally including or excluding nodes within the subtree.

The optimal $\beta_j$ can be found the same way as on the bottom level but this time we need to reorder $D_{jk}=\frac{1}{2}\|m|_{T_{jk}}\|^2-F_{jk}-C_j$, the difference in the weight of including or pruning a branch. The algorithm is summarised in Algorithm \ref{alg:beta}. 

Notice that we have selected $\hat{\boldsymbol{\beta}}$ recursively from bottom up. It is possible that a node is considered to contain enough information on level $j_1$, leading to a certain $\hat{\beta}_{j_1}$, but the whole branch is later pruned on a level $j_2<j_1$. This would allow changing the value of $\hat{\beta}_{j_1}$ without affecting the final tree.


\begin{algorithm}[ht!]
\caption{Pseudo-code for automatic $\beta$ selection with Gaussian base prior}
\label{alg:beta}
$j=J$\\
  \For{$k = 0$ \KwTo $2^J - 1$}
  {Set $F_{Jk}=\frac{1}{4}m^2_{Jk}$\\
    Calculate grid points $\beta_{Jk} = \big(1+\exp(\frac{1}{4}m^2_{Jk})\big)^{-1}$
    \\
    Denote the set of included nodes $A_{\beta_{Jk}}=\{(J,\tilde{k}) \g \beta_{J\tilde{k}}\geq\beta_{Jk}\}$
    }
    
Find $\hat{\beta}_{J} =\arg\min \big(L(m_J,\beta_{Jk})+R(\beta_{Jk})\big)$ where\\
$L(m_J,\beta_{Jk})=\sum_{k\in A_{\beta_{Jk}}}F_{Jk}+\sum_{k\notin A_{\beta_{Jk}}}\frac{1}{2}m^2_{Jk}$ \\
$R(\beta_{Jk})=|A_{\beta_{Jk}}|\log\Big(\frac{1-\beta_{Jk}}{\beta_{Jk}}\Big)
-2^J\log\bigg((1-\beta_{Jk})\Big(\frac{1}{2}-\beta_{Jk}\Big)^a\bigg)$
\vspace{4mm}

$j = j-1$\\
\While{$1 \leq j \leq J$}
    {\For{$k = 0$ \KwTo $2^j - 1$}
    {For the left child $(j',k')$ of node $(j,k)$\\
    \If{
    $F_{j'k'}-\log(\hat{\beta}_{j'})
    <\frac{1}{2}\|m|_{T_{j'k'}}\|^2-(2^{J-j}-1)\log(1-\hat{\beta}_{j'})$}
   {$t_{j'k'}=1$ and 
   $F_{jk}=\frac{1}{4}m_{jk}^2+F_{j'k'}-\log(\hat{\beta}_{j'})$}
   \Else{
   $t_{j'k'}=0$ and 
   $F_{jk}=\frac{1}{4}m^2_{jk}+\frac{1}{2}\|m|_{T_{j'k'}}\|^2-\sum_{\ell=j'}^J 2^{\ell-j'}\log(1-\hat{\beta}_{\ell})$}
   
For the right child $(j'',k'')$ of node $(j,k)$\\
   \If{
   $F_{j''k''}-\log(\hat{\beta}_{j''})
    <\frac{1}{2}\|m|_{T_{j''k''}}\|^2-(2^{J-j}-1)\log(1-\hat{\beta}_{j''})$}
   {$t_{j''k''}=1$ and 
   $F_{jk}=F_{jk}+F_{j''k''}-\log(\hat{\beta}_{j''})$}
   \Else{
   $t_{j''k''}=0$ and 
   $F_{jk}=F_{jk}+\frac{1}{2}\|m|_{T_{j''k''}}\|^2-\sum_{\ell=j''}^J 2^{\ell-j''}\log(1-\hat{\beta}_{\ell})$}
   
   Calculate grid points $\beta_{jk} = \big(1+\exp(\frac{1}{2}\|m|_{T_{jk}}\|^2-F_{jk}-C_j)\big)^{-1}$ where\\ 
   $C_j=\sum_{\ell=j+1}^J 2^{\ell-j} \log((1-\hat{\beta}_\ell)(\frac{1}{2}-\hat{\beta}_\ell)^a)$\\
   Denote the set of included nodes 
   $A_{\beta_{jk}}=\{(j,\tilde{k}) \g \beta_{j\tilde{k}}\geq\beta_{jk}\}$
   }
   Find $\hat{\beta}_{j} =\arg\min \big(L(m_j,\beta_{jk})+R(\beta_{jk})\big)$ where\\
    $L(m_j,\beta_j)= \sum_{k\in A_{\beta_j}}F_{jk} 
    + \sum_{k\notin A_{\beta_j}}\Big(\frac{1}{2}\|m|_{T_{jk}}\|^2-C_j\Big)$ \;
    $R(\beta_{jk})=|A_{\beta_{jk}}|\log\Big(\frac{1-\beta_{jk}}{\beta_{jk}}\Big)
    -2^j\log\bigg((1-\beta_{jk})\Big(\frac{1}{2}-\beta_{jk}\Big)^a\bigg)$
    } 

$\tilde{t}_{00}=1$\\
\While{$1 \leq j \leq J$}{
  \For{$k = 0$ \KwTo $2^j - 1$}{
  $\wt_{jk}=\prod_{(j',k')\succeq(j,k)}t_{j'k'} $\\
    \If{$\tilde{t}_{jk} = 1$}{
      $g_{jk} = m_{jk}$
    }
    \Else{
      $g_{jk} = 0$
    }
  }
  $j = j-1$
}
\vspace{8mm}
\end{algorithm}



\begin{proposition}\label{Prop:scaling}
Consider the denoising problem of estimating $f$ from a noisy measurement \eqref{eq:DenoisingMain}. We assume a truncated random tree Besov prior for $f$ as described in \eqref{eq:DefPrior}. Denote by $T^{1}_\beta(m_1)$ the denoised signal arising from tree pruning with a fixed $\beta$ when $w_i\sim \Nn(0,1)$ and $g_{jk}\sim\Nn(0,1)$. 

Then, for $\sigma,\kappa>0$, the denoised signal $T^{\kappa}_\beta(m_\sigma)$ arising from a measurement with noise $w_i\sim\Nn(0,\sigma^2)$ and prior choice $g_{jk}\sim\Nn(0,\kappa^2)$ is given by $\tilde{T}^1_{\tilde{\beta}}(m_1)$, where the coefficients $\hat{g}^1_{jk}=m_{jk}/2$ are replaced by $\hat{g}^1_{jk}=\kappa^2m_{jk}/(\kappa^2+\sigma^2)$ and
\begin{align*}
    \widetilde{\beta}=\frac{\beta^{c}}{\beta^{c}+(1-\beta)^{c}}
\end{align*}
with $c=\sigma^2(\kappa^2+\sigma^2)/(2\kappa^2)>0$.
\end{proposition}

Scaling the coefficients $\hat{g}^1_{jk}$ means scaling the denoised signal, which is undesirable. For this reason, in practice we set $\hat{g}^1_{jk}=m_{jk}$. Proposition \ref{Prop:scaling} then shows that once the optimal $\hat{\beta}^1_1$ has been estimated for $T^1_\beta(m_1)$ the corresponding optimal value $\hat{\beta}^\kappa_\sigma$ for $T^{\kappa}_\beta(m_\sigma)$ can be obtained directly from $\hat{\beta}^1_1$. 
Equivalently, one may view the noise and prior variances, 
$\sigma^2$ and $\kappa^2$, as influencing only the optimal choice of the regularisation parameter $\beta$.

The following remark shows how scaling the input signal influences the resulting denoised signal.

\begin{remark}
We see from Proposition \ref{Prop:scaling} that, for a constant $c$, $T^1_\beta(cm_1) =cT^1_{\widetilde{\beta}}(m_1)$ where 
\begin{align*}
    \widetilde{\beta}=\frac{\beta^{c^2}}{\beta^{c^2}+(1-\beta)^{c^2}}.
\end{align*}
\end{remark}

\begin{proof}[Proof of Proposition \ref{Prop:scaling}]
We start by noting that when $\beta$ is fixed and $w_{jk}\sim\Nn(0,\sigma^2)$, $g_{jk}\sim\Nn(0,\kappa^2)$ the posterior of $g_{jk}\g \tilde{t}_{jk}=1$ is given by
\begin{align*}
z_{jk}^1 & \propto \exp\Big(-\frac{1}{2\sigma^2}(m_{jk}-g_{jk})^2
-\frac{1}{2\kappa^2}g_{jk}^2)\Big)\beta
\end{align*}
and the posterior for $g_{jk}\g \tilde{t}_{jk}=0$ is 
\begin{align*}
z_{jk}^0 & \propto \exp\Big(-\frac{1}{2\sigma^2}m_{jk}^2-\frac{1}{2\kappa^2}g_{jk}^2)\Big)(1-\beta).
\end{align*}

In the bottom row $j=J$ the weight resulting from including a node $(J,k)$ is simply
\begin{align*}
    Y_{Jk}^1=\min_{g_{Jk}}\Big(-\log(z_{Jk}^1)\Big)
=\frac{m_{Jk}^2}{2(\kappa^2+\sigma^2)}-\log(\beta),
\end{align*}
which is achieved at $\hat{g}_{Jk}=\kappa^2m_{Jk}/(\kappa^2+\sigma^2)$, and the weight for not including the node is 
\begin{align*}
    Y_{Jk}^0=\min_{g_{Jk}}\Big(-\log(z_{Jk}^0)\Big)
=\frac{m_{Jk}^2}{2\sigma^2}-\log(1-\beta),
\end{align*}
following from the choice $\hat{g}_{Jk}=0$.

A node $(J,k)$ is included in the subtree if including the node carries less weight than turning it off, that is, when $Y_{Jk}^1\leq Y_{Jk}^0$ or equivalently when 
\begin{align*}
    \log\Big(\frac{1-\beta}{\beta}\Big)
    \leq \frac{\kappa^2m_{Jk}^2}{2\sigma^2(\kappa^2+\sigma^2)}
    =\frac{1}{c}\frac{m_{Jk}^2}{4},
\end{align*}
where $c=\sigma^2(\kappa^2+\sigma^2)/(2\kappa^2)>0$.  We can rewrite the above as 
\begin{align*}
    \log\Bigg(\Big(\frac{1-\beta}{\beta}\Big)^c\Bigg)
    =\frac{m_{Jk}^2}{4},
\end{align*}
which corresponds to 
\begin{align*}
    \log\Big(\frac{1-\widetilde{\beta}}{\widetilde{\beta}}\Big)
    =\frac{m_{Jk}^2}{4},
\end{align*}
where $\widetilde{\beta}=\beta^c/(\beta^c+(1-\beta)^c)\in[0,0.5]$ for $\beta\in[0,0.5]$ and $c>0$. 
\end{proof}


\subsection{Automatic $\beta$ selection with Laplace draws}\label{Sec:PruningLaplace}

Optimising $\boldsymbol{\beta}$ when $g_{jk}\sim \text{Laplace}(0,\kappa)$ is similar to the Gaussian case though we need to take into account the additional soft thresholding which forces $\hat{g}_{jk}=0$ for all $|m_{jk}|\leq\frac{1}{\kappa}$. On the bottom level $j=J$ the weight of including a node is 
\begin{align*}
Y^1_{Jk} 
= \min_{g_{Jk}}\big(-\log(z^1_{Jk})\big)
=\begin{cases}
\frac{1}{2}m_{Jk}^2 -\log\Big(\beta_J\big(\frac{1}{2}-\beta_J\big)^a\Big) 
& \text{when } |m_{Jk}|\leq \frac{1}{\kappa} \\
\frac{1}{\kappa}|m_{Jk}|-\frac{1}{2\kappa^2}
-\log\Big(\beta_J\big(\frac{1}{2}-\beta_J\big)^a\Big)  
& \text{when } |m_{Jk}|> \frac{1}{\kappa},
\end{cases}
\end{align*}
which is attained at 
\begin{align*}
\hat{g}^1_{Jk}=
\argmin_{g_{Jk}}\big(-\log(z^1_{Jk})\big)
=\begin{cases}
      0 & \text{when}\ |m_{Jk}|\leq \frac{1}{\kappa} \\
     m_{Jk}-\text{sign}(m_{Jk})\frac{1}{\kappa} & \text{when}\ |m_{Jk}|> \frac{1}{\kappa}.
 \end{cases}
\end{align*}
As in the Gaussian case the weight of a pruned node is 
\begin{align*}
Y^0_{Jk} 
= \min_{g_{Jk}}\big(-\log(z^0_{Jk})\big)
=\frac{1}{2}m_{Jk}^2 -\log\Big((1-\beta_J)\big(\frac{1}{2}-\beta_J\big)^a\Big)
\end{align*}
which follows from setting $\hat{g}^0_{Jk}=0$.

A bottom row node is included in the tree if $Y^1_{Jk}\leq Y^0_{Jk}$. We immediately see that a node that carries little information $|m_{Jk}|<1/\kappa$ is always pruned on the bottom level since the penalty for including it is larger and there is no gain on the data fit. 

A bottom node is included in the tree if $|m_{Jk}|>1/\kappa$ and 
\begin{align}
\frac{1}{\kappa}|m_{Jk}|-\frac{1}{2\kappa^2}-\log(\beta_J)
\leq \frac{1}{2}m^2_{Jk}-\log(1-\beta_J).
\end{align}
In particular, given a measurement $|m_{Jk}|>1/\kappa$ the node $(J,k)$ is selected if 
\begin{align*}
\beta_J\geq\frac{1}{1+\exp\big(\frac{1}{2}m^2_{Jk}-\frac{1}{\kappa}|m_{Jk}|+\frac{1}{2\kappa^2}\big)}.    
\end{align*}

As before we want to solve 
\begin{align*}
\hat{\beta}_J = \argmin_{\beta_{Jk}} B_{J,A_{\beta_{Jk}}} 
= \argmin_{\beta_J} \big(L(m_J,\beta_{Jk})+R(\beta_{Jk})\big),
\end{align*}
where the penalty term $R(\beta_J)$ is the same as in \eqref{eq:WeightBottomGaussian} and the data term is
\begin{align*}
L(m_J,\beta_{Jk}) = \sum_{k\in A_{\beta_{Jk}}}\bigg(\frac{1}{\kappa}|m_{Jk}|-\frac{1}{2\kappa^2}\bigg)
+\sum_{k\notin A_{\beta_{Jk}}} \frac{1}{2}m^2_{Jk}.
\end{align*}
The minimiser can be found similarly to the Gaussian case, but it suffices to consider only the grid points 
$\beta_{Jk}=(1+\exp(m^2_{Jk}/2-|m_{Jk}|/\kappa+1/{2\kappa^2}))^{-1}$ corresponding to nodes $(J,k)$ where $|m_{Jk}|>1/\kappa$.  

Once $\hat{\beta}_J$ has been found we can move to the next level. When $j<J$ including a node into the tree can be beneficial even when $|m_{jk}|<1/\kappa$ if the optimised subtree starting from the node $(j,k)$ carries enough information. In this case we set $\hat{g}^1_{jk}=0$, $\tilde{t}_{jk}=1$, which provides additional regularisation within the optimised subtree. 



\section{Numerical Results \label{Sec:Numerical}}
We will now showcase the performance of the proposed automatic 
$\beta$-selection tree pruning method in 1D and 2D denoising, and 1D deconvolution.  
%The 2D experiments are conducted on open-access 2D images available at the Kodak Lossless True Color Image Suite. 


\subsection{Signal denoising and deconvolution}

In our first example, we consider the blocks test data from \cite{Donoho1994}, displayed on the left in Figure \ref{Denoise1D}. We assume that only a noisy signal, with signal to noise ratio $5$, is observed. We denoise the signal using random tree Besov priors with Haar wavelet basis. 

\begin{figure}[t]
\centering
\includegraphics[width=\textwidth]{Denoise1D.png}
    \caption{The original signal in black and the  measurement in red (left), denoised signal using Gaussian base prior and automatic $\beta$ selection (middle), and denoised signal using Laplace base prior and automatic $\beta$ selection (right).}
    \label{Denoise1D}
\end{figure}

For the automatic $\beta$ selection method we used hyperparameter $a=100$ but the algorithm is not sensitive to the exact value of $a$. Values between $15$ and $500$ produce the exact same reconstruction and values between $5$ and $1000$ produce very similar reconstructions with only a small difference in the relative error. 

The denoised signal under Gaussian base prior with automatic $\beta$ selection is shown in the middle in Figure \ref{Denoise1D}. The denoised signal corresponding to optimally chosen fixed $\beta=10^{-4}$ is visually indistinguishable. The relative errors between the original signal and the denoised signal arising from the automatic $\beta$ selection method and from the optimal fixed $\beta$ method are $1.116\%$ and $1.110\%$, respectively. 

\begin{figure}[b]
\centering
\includegraphics[width=\textwidth]{GaussianTreeBeta.png}
    \caption{ Selected wavelet coefficient tree for the denoised signal under the Gaussian base prior (left) and the levelwise estimates 
    $\log(\hat{\boldsymbol{\beta}})$ (right).}
    \label{GaussianTreeBeta}
\end{figure}

The selected wavelet coefficient tree for the denoised signal under the Gaussian base prior with automatic $\beta$ selection is shown in Figure \ref{GaussianTreeBeta} on the left and the values $\log(\hat{\beta}_j)$ for $j=3,\dots,13$ are plotted on the right. We omit $j=1,2$ because $\hat{\beta}_1$ and $\hat{\beta}_2$ are extremely small and would dominate the scale. At the coarsest levels the cost difference between keeping a branch and pruning it is very large, so the data-determined $\beta$-grid concentrates near zero and the optimisation selects all nodes at those levels. From level $j=3$ onwards, pruning becomes non-trivial and some branches are removed, as can be seen in the tree plot. At the finest scales most coefficients carry little information, leading to extensive pruning; correspondingly, the estimated $\hat{\beta}_j$ at fine levels are small.

The denoised signal under the Laplace base prior with automatic $\beta$ selection is shown on the right in Figure \ref{Denoise1D}. Here too, the denoised signal corresponding to the optimally tuned fixed $\beta=10^{-4}$ is visually indistinguishable. The relative errors between the original signal and the denoised signal obtained with the automatic $\beta$ selection method and with the optimally tuned fixed-$\beta$ method are $1.185\%$ and $1.150\%$, respectively. We used $\kappa=1$ for both methods and $a=10$ for the automatic $\beta$ selection. The method is somewhat more sensitive to the choice of the hyperparameter $a$ than in the Gaussian base prior case.

We see that automatic $\beta$ selection yields near-optimal performance while eliminating manual tuning of the regularisation parameter. The denoised signals obtained under the Gaussian and Laplace base priors are very similar. In these experiments, the choice of wavelet basis appears to have a larger effect on the reconstruction than the choice of base prior.

As a second 1D example we consider a mildly ill-posed deconvolution problem. The observed signal is given by
\begin{align}
M=Af+W,
\end{align}
where $A$ is a known convolution operator, $f$ is a piecewise linear function, and the measurement contains $1\%$ Gaussian noise. The original and blurred noisy signals are shown on the left in Figure \ref{Deconvolution}. We solve the inverse problem using a plug-and-play scheme \cite{Venkatakrishnan2013}. In each iteration, the linear data-fidelity step for the convolution model is followed by a denoising step that applies our random tree Besov prior with Gaussian base prior and automatic $\beta$ selection to the current iterate.

The deconvolved signal is shown on the right in Figure \ref{Deconvolution}. The reconstruction exhibits a characteristic staircasing behaviour, similar to solutions obtained with total variation priors, while still capturing the main discontinuities accurately. This example is intended as a proof of concept demonstrating that the proposed automatic $\beta$-selection tree prior can be used as a plug-and-play denoiser inside more general inverse problem solvers; a more systematic study of deconvolution is left for future work.
 

\begin{figure}[t]
\centering
\includegraphics[width=\textwidth]{Deconvolution.png}
    \caption{The original signal in black, measurement in blue and deconvolved signal in red.}
    \label{Deconvolution}
\end{figure}

\subsection{Image denoising}

The 2D experiments are conducted on open-access 2D images available at the Kodak Lossless True Color Image Suite. We compare the MAP estimators arising from random tree Besov priors with automatic $\beta$ selection and fixed optimally tuned $\beta$. The optimal regularisation parameter $\beta$ is chosen to maximise the Structural Similarity Index (SSIM) between the ground truth image and the denoised one. The original image and the noisy measurement, containing $7\%$ Gaussian noise, are shown in Figure \ref{Data2D}.   

\begin{figure}[h]
\centering
\includegraphics[width=\textwidth]{OriginalNoisy.pdf}
    \caption{The original image and the noisy measurement with $7\%$ Gaussian noise.}
    \label{Data2D}
\end{figure}

We use Daubechies 2 wavelets to decompose the images, as they mitigate blocky artefacts while preserving edges through smooth transitions. The image dimension is $512\times 512$ and the original image values are in the interval $[0,1]$. For tree pruning with the automatic $
\beta$ selection method we scale the values to match the strength of the hyperprior. Based on testing, we established a heuristic rule that the scale value should be in $[200/\delta,300/\delta]$, where $\delta$ is the percentage of noise. 

For the fixed $\beta$ pruning method with Gaussian base prior, the optimal regularisation parameter, based on the SSIM, is calculated to be $\beta=0.4986$ and the SSIM  is equal to $0.7992$. 
With Laplace base prior the threshold parameter is set to $\kappa=0.11.$ For the soft pruning method with fixed $\beta$, the optimal regularisation parameter is $\beta=0.4993$ and the SSIM  is equal to $ 0.7715$. Note that $\kappa$ should be optimised separately and the optimal value will be different for fixed and automatic $\beta$ selection methods.
    
With Gaussian base prior pruning with automatic $\beta$ selection results in SSIM equal to $0.7967$. With Laplace base prior soft pruning with automatic $\beta$ selection results in SSIM equal to $0.7852$. As mentioned above the latter could be improved by optimising $\kappa$ separately.

All of the tree pruning methods outperform optimal soft and hard thresholding applied to the wavelet coefficients. See Table \ref{tab:denoising_results L2} for the results.

\begin{figure}[ht]
\centering
\includegraphics[width=\textwidth]{FoGResults.pdf}
    \caption{Denoised images using Gaussian base prior and automatic $\beta$ selection (top left), Gaussian base prior and optimal fixed $\beta$ (top right), Laplace base prior and automatic $\beta$ selection (bottom left), and Laplace base prior and optimal fixed $\beta$.}
    \vspace{0.5cm}
    \label{Denoised2D}
\end{figure}


\begin{table}[ht]
\vspace{0.5cm}
    \centering
    \caption{Comparison of tree pruning methods based on SNR, SSIM, and Relative Error for the Gaussian and Laplace base priors.}
    \label{tab:denoising_results L2}
    \begin{tabular}{|l|c|c|c|}
        \hline
        \textbf{} & \textbf{SNR (dB)} & \textbf{SSIM} & \textbf{Relative Error} \\
        \hline
        Noisy Image & 18.0382 & 0.4238 & 0.2758 \\
        Tree pruning-fixed $\beta$ (Gaussian) & 18.6696 & 0.7992 & 0.2585  \\
        Tree pruning-automatic $\beta$ selection & 18.6577 & 0.7967 & 0.2589 \\
        Soft tree pruning-fixed $\beta$ (Laplace) & 18.4024 & 0.7715 & 0.2683  \\
        Soft tree pruning-automatic $\beta$ selection & 18.5162 & 0.7852 & 0.2634 \\
         Soft thresholding  & 18.4074 & 0.7704 & 0.2682 \\
        Hard thresholding  & 18.5728 & 0.7740 & 0.2617\\
        \hline
    \end{tabular}
\end{table}


 \newpage
%=====================================

\vspace{1cm}

\bibliographystyle{siam}
\bibliography{ref}
%
\end{document}
